\section{时间外划分与同信息强对照}
\subsection{扩展窗口和结果区间剔除}
测试折依次为2021、2022、2023和2024年截至7月24日，验证使用测试前一年，训练窗口逐年扩展。图\ref{fig:folds}只按年格呈现结构，实际索引按市场日和结果结束日执行边界检查。训练标签结果结束日必须严格早于验证首个市场日，验证标签结果结束日必须严格早于测试首个市场日。

\reportfigure{四段扩展训练、独立验证与时间外测试}{fig:folds}{\FigFolds}{已生成的四折共享索引设计。年格不表示精确日长度，2024只覆盖至7月24日，purge依据结果结束日。}

公司日索引只有一份，模型通过信息权限读取不同输入。训练序列可以向前读取验证日前的历史，不能向后读取结果区间；关系邻居使用预测日已知的历史表示。缺失公告或关系时，保留目标标签和序列，用掩码表达辅助不可用。另报共同覆盖子样本有助于理解关系稀疏程度，但主要比较仍使用共同全量目标。

时间序列交叉验证的有效性依赖具体条件\cite{r33}，不能由某些条件下可用随机划分推出股票公司日也可以打散。市场共同冲击、重复公告家族和五日标签重叠同时存在，独立年份与结果区间剔除更贴近这里的预测问题。预处理、文本降维、校准与选模各自限定拟合集合，避免模型在测试分布上调整后再宣称时间外表现。

\subsection{比较组和输入预算}
收益-only组包括下行波动/传统尾部统计Logit、强树和GRU。历史特征可含滚动残差偏度、上下行波动比、波动和回撤，全部由起点已知资料计算。LightGBM、XGBoost、CatBoost是成熟非线性表格工具\cite{r25,r26,r27}，正式比较应给予合理验证预算，而不是仅使用默认参数的弱线性模型。

同信息组取得原生及预测事件、正文聚合、诉讼、关系统计、边龄和邻居历史量。强树取得同一正文向量的训练期PCA64表示，不能只给它公告数量，却让主模型独享语义信息。同编码器组则保持同GRU、事件缓存和二部消息，主要改变拼接、固定融合、普通门控、MSGCA式交叉注意力适配和本方法。

\reporttable{比较模型的信息权限及适配说明}{tab:baselines}{
\begin{tabularx}{\linewidth}{p{3.05cm}YY}\toprule
模型组&信息权限&比较目的\\\midrule
收益统计Logit&收益、三因子及训练期历史统计&量化简单风险特征能力\\
收益强树/GRU&同收益主信息&检验序列表示与非线性表格差别\\
同信息强树&同正文、事件、诉讼和图统计&排除信息剥夺造成的虚假增量\\
同编码器融合&同三路表示和结构控制&检验融合约束，而非换编码器\\
图时序适配&同标签、日期与允许关系&研究直接图方向在本任务下的表现\\\bottomrule
\end{tabularx}}

Graph WaveNet及多关系风险方法按同一五日标签适配。直接近邻完整方法或代码尚未取得时，只能标为启发式适配，不能称忠实复现。每折主对照从全部同信息模型中按验证种子平均AP选择，排除本研究模型M；测试不参与选择。树模型每信息组每树最多8个验证配置，神经模型采用共同候选和停止规则，保留全部选参记录。

\section{风险排序、概率质量和配对统计}
\subsection{指标的分工}
AP用于衡量正类风险排序。Davis和Goadrich以及Saito和Rehmsmeier讨论了PR与ROC在不平衡评价中的关系\cite{r28,r29}。本数据P10正类占比约10.94\%，仅凭多数类就能得到较高分类准确率，因而准确率无法充分说明有限审核名单的价值。AP也受正类率影响，需逐折报告实际正类率、AP和差值。

每日前10\%取$k_t=\lceil0.1N_t\rceil$，分数同分时按证券代码固定排序。报告每日精确率及有正例日期的召回均值，而不是把全年公司日混排后选前10\%。F1阈值和温度校准只在验证集拟合；ROC-AUC和Brier作为辅助评价，与主排序增量分开解释。模型概率在未经验证校准以前只称风险分数。

\subsection{种子与日期不确定性分别计算}
双方每折固定种子42、2024、7、123、3407。令$AP_{f,s}^{M}$和$AP_{f,s}^{B}$为模型与该折验证选出主对照的指标，主差值为
\begin{equation}
 \Delta AP=\frac14\sum_{f=1}^{4}\left(\frac15\sum_{s=1}^{5}AP_{f,s}^{M}-\frac15\sum_{s=1}^{5}AP_{f,s}^{B}\right).
\end{equation}
先计算各模型各个种子AP，再取均值。AP对排序是非线性的，所以$AP(\frac15\sum_sp_s)$通常不等于$\frac15\sum_sAP(p_s)$；概率集成可以另报，但不是同一个估计对象。

\reportfigure{同公司日配对、种子均值和日期块区间的计算路径}{fig:evaluation}{\FigEvaluation}{统计方案，不包含模拟预测分数或区间。日期块在折内抽取，再进行四折等权汇总。}

K\"unsch和Politis、Romano的工作为依赖序列重采样提供了基础\cite{r31,r32}。本课题采用整市场交易日的配对移动块，20日为主块长，5/60日为敏感性，1000次。一次抽取中，同一天的全部公司及两模型预测一起进入样本，保留横截面共同冲击。这里不是平稳bootstrap的随机长度算法原样复现，具体抽样规则须与实现记录一致。

训练种子标准差刻画优化不确定性，日期块区间刻画样本时序不确定性，两者分别报告。四折不按样本量加权，以免公司日更多的年份支配结论。主对照选择发生在验证期，区间计算发生在冻结测试预测后；多次试探和事后选择会影响解释，White关于数据窥探的讨论对此有直接警示\cite{r34}。

\section{十一项消融与机制压力}
消融区分三个问题：信息是否存在，限制是否有效，结果是否来自不可靠来源或关系巧合。表\ref{tab:ablations}中的11项保持共同样本和训练协议；P5作为稳健性标签，普通门控和拼接属于架构对照，不混入这11项计数。

\begin{longtable}{p{.65cm}p{3.1cm}p{4.5cm}p{5.3cm}}
\caption{固定消融、改变内容和所检验问题}\label{tab:ablations}\\
\toprule 序号&实验&改变内容&检验问题\\\midrule\endfirsthead
\toprule 序号&实验&改变内容&检验问题\\\midrule\endhead
1&去公告&关闭预告分支&公告信息是否提供额外效用\\
2&去图&关闭关系分支&关系系统整体增量\\
3&去全部年龄&关闭事件注意力、边消息及修正年龄&年龄处理整体作用\\
4&去质量&移除质量过滤/加权&接口质量策略是否有效\\
5&年龄仅作特征&保留年龄输入，无幅度包络&信息变量与幅度限制差别\\
6&去辅助损失&只用主任务BCE&分支监督与融合约束的作用\\
7&仅去边级衰减&保留图外包络&相对邻居权重作用\\
8&度保持随机图&随机重连关系&真实关系与容量/度数效应\\
9&仅原生公告&不使用Qwen语义&语言适配的必要性\\
10&仅Qwen结构&移除正文隐藏表示&结构抽取与连续语义差别\\
11&匿名文本&匿名公司/证券包装&预训练记忆风险的敏感性\\\bottomrule
\end{longtable}

去质量控制可能同时改变有效事件集合，不能把差值解释为纯粹标量门效应。随机图应保持度数约束、预测日可见性和机构独享规则，不把重连引入的非法未来边当作关系随机性。匿名化仅减轻记忆线索，无法证明模型未记住事件内容。每项实验需将变更的数据产品、参数和输入权限列入记录。

自然分组按模型输入中可获得的年龄、质量和共享邻居规模划分，观察新旧事件、低覆盖公司和不同市场年份。人为增龄保持内容与权重固定，检查允许上界与实际分数变化。这种压力检验用于解释机制，不作为新测试样本或额外金融收益结论。若总体增量仅由某一年或高覆盖群体驱动，需明确适用条件。

\section{日装配、存储与训练恢复设计}
\subsection{共享数据而不是重复宽表}
把6,116,554个标签全部展开成60日、五字段float32窗口，仅数值本身就约为
\begin{equation}
 6{,}116{,}554\times60\times5\times4\approx7.34\times10^9\text{ 字节}.
\end{equation}
它不包含掩码、索引、正文向量、图或运行时副本，已接近当前主机内存尺度。研究拟保存原始序列与位置索引，通过市场日查找窗口；事件缓存全局共享，不按公司日重复保存。上述数值是存储理论下限，不是已经测得的程序峰值。

日装配器查询目标和邻居的并集，统一标准化后返回序列、事件、边及监督掩码。当天模型计算一次公司表示，再用稀疏索引聚合机构。训练目标集合和计算节点集合必须分开；某公司作为邻居出现，不自动获得未来标签监督。跨日期合并若改变图、年龄或控制变量，则不能简化为一个静态全样本图。

\reporttable{数据产品及版本失效条件}{tab:storage}{
\begin{tabularx}{\linewidth}{p{3.1cm}YY}\toprule
产品&保存内容&何时失效或需重建\\\midrule
标签及四折索引&日期、证券、结果结束日、标签、split&收益/因子、标签算法或范围变化\\
原始序列索引&市场位置、字段及掩码&字段顺序、日历、标准化版本改变\\
正文缓存&原文哈希、模型/模板、正文位置、预测字段&任一缓存键组成变化\\
完整图版本流&公司、主体、比例、公开版本、生效日&名单/主体/证券对应规则修订\\
批次manifest&输入版本、维度、允许信息、目标集合&任何产品或装配合同变化\\
冻结预测&模型、索引、种子、日期与输出&模型/输入改变，不覆盖旧结果\\\bottomrule
\end{tabularx}}

\subsection{复杂性与计算瓶颈}
设$N$为单日去重公司数、$J$为机构数、$E$为边数，序列长度$T=60$，实际GRU输入通道$F^*=10$，隐藏维$h=64$，投影$d=128$。序列计算约为$O(TN(F^*h+h^2))$；二部变换及聚合约为$O((N+J)d^2+Ed)$；事件聚合受$K\le8$限制。编码当前公司一次后再共享给各机构，避免重复GRU成本。

大模型正文编码成本发生在事件级缓存阶段，不能隐藏在风险网络每步耗时之外。资源日志分别记录抽取、缓存、装配、风险训练和评价I/O。高度机构采用分块全邻居聚合；采样策略先在小批次比较输出偏差和截断率，不能只报告节省内存的结果。实际小时数、显存和费用在GPU试跑后测量，不用合成CPU步时外推。

\subsection{恢复运行与异常处理}
检查点绑定模型配置、优化器、随机种子、输入版本和训练位置。恢复后不得重新选择验证表现更好的种子或合并不同缓存；预测按共同索引检查重复、漏行和年份边界。记录可重复命令和环境，特别区分CPU与GPU依赖、量化库及模型revision。

\reporttable{常见失败、影响与定位方法}{tab:failure}{
\begin{tabularx}{\linewidth}{p{3.15cm}YY}\toprule
失败情况&可能影响&定位或处理\\\midrule
代码命名空间混用&公告/边连到错误证券&源字段追溯，唯一历史区间确认\\
名单口径混合&重复边及假共享关系&类型分离、原版比对，不按排名删行\\
历史覆写不可追溯&日期正确但内容有未来信息&版本政策与原公告，未决资料隔离\\
正文范围或截断变化&缓存包含答案或与模型不一致&token位置检查及缓存键失效\\
真实批次内无监督目标&错误训练或空损失&立即失败，不用邻居标签补目标\\
恢复后漏预测/重复索引&配对评价样本改变&全量键检查，按冻结索引重读\\
年龄约束过强&有效旧信息被削弱&普通门控、自然分组和增龄压力比较\\\bottomrule
\end{tabularx}}

异常处理应使问题可见。原版匹配失败、缺少ID、解析失败或极小图分母都需要独立计数，不能用一个“缺模态”标签掩盖全部原因。已完成输入检查并不意味着标签或主体永远正确，后续修订必须保持旧数据版本和结果，重新确认受影响实验。
